% Môn: Vật lý
% Lớp: 12
% Chương: Chương I. Vật lí nhiệt
% Bài: L12C1 Bài 4. Nhiệt dung riêng · Bài toán tính toán nhiệt lượng cần thiết để hóa hơi hoàn toàn một lượng chất lỏng ở nhiệt độ sôi (áp dụng công th�[...]
% ===== Câu 44 =====
% ID: L12C1B4-14-TN
% Mức: VD
\dangbt{Bài toán tính toán nhiệt lượng cần thiết để hóa hơi hoàn toàn một lượng chất lỏng ở nhiệt độ sôi (áp dụng công thức Q =}
\begin{ex}
% ID: L12C1B4-14-TN
% Mức: VD
Biết nhiệt độ sôi, nhiệt dung riêng và nhiệt hóa hơi riêng của nước lần lượt là $100^\circ C$, $4200\ J/kg.K$ và $2,3.10^6\ J/kg$. Nhiệt lượng cần cung cấp để [...]
\choice
{$2,6.10^6\ J$}
{\True $5,3.10^6\ J$}
{$26.10^6\ J$}
{$53.10^6\ J$}
\loigiai{
Nhiệt lượng cần cung cấp để làm hóa hơi hoàn toàn $2\,\text{kg}$ nước từ nhiệt độ $20^{\circ} {C}$ là:
${Q}={mc} \Delta {t}+{mL}=2 . 4200 .(100-20)+2 . 2,3 . 10^{6}=5272000 {~J} \approx 5,3 . 10^{6} {~J}$.
}
\end{ex}
% =========================================================
% BÀI 1
% =========================================================
\begin{ex}
Một xô có chứa $M=10\,\mathrm{kg}$ hỗn hợp nước và nước đá ở
trong phòng. Sự thay đổi nhiệt độ của hỗn hợp theo thời gian được
biểu diễn như hình vẽ. Lấy gần đúng nhiệt dung riêng của nước là
$c=4200\,\mathrm{J/(kg.K)}$, nhiệt nóng chảy riêng của nước đá là
$\lambda=3,4\cdot10^5\,\mathrm{J/kg}$. Cho rằng sự hấp thụ nhiệt từ
môi trường là đều.

\begin{center}
\begin{tikzpicture}[>=stealth,scale=0.9]
    \draw[->] (0,0) -- (0,3) node[left] {$t(^\circ\mathrm C)$};
    \draw[->] (0,0) -- (5,0) node[below] {$\tau(\text{phút})$};

    \draw[dashed] (0,2.5) -- (4.2,2.5);
    \draw[dashed] (3.5,0) -- (3.5,2.5);
    \draw[dashed] (4.2,0) -- (4.2,2.5);

    \draw[thick] (3.5,0) -- (4.2,2.5);

    \node[left] at (0,2.5) {$2$};
    \node[below] at (3.5,0) {$50$};
    \node[below] at (4.2,0) {$60$};

    \fill (3.5,0) circle (1.5pt);
    \fill (4.2,2.5) circle (1.5pt);

    \node[below left] at (3.5,0) {$A$};
    \node[above right] at (4.2,2.5) {$B$};
\end{tikzpicture}
\end{center}


Xét tính đúng -- sai của các phát biểu sau:

\choiceTF
{\True Tại điểm $A$, toàn bộ nước đá trong xô đã tan hết.}
{Trong $50$ phút đầu, hỗn hợp không hấp thụ nhiệt từ môi trường.}
{\True Khối lượng nước ban đầu trong xô xấp xỉ $8,76\,\mathrm{kg}$.}
{\True Khối lượng nước đá còn lại trong xô ở thời điểm $20$ phút xấp xỉ
$0,75\,\mathrm{kg}$.}

\loigiai{
Từ phút thứ $50$ đến phút thứ $60$, nhiệt độ nước tăng từ
$0^\circ\mathrm C$ lên $2^\circ\mathrm C$.

Nhiệt lượng nước nhận trong $10$ phút cuối:
\[
Q=Mc\Delta t
=10\cdot4200\cdot2
=84000\,\mathrm J.
\]

Do nhiệt lượng cung cấp đều nên trong $50$ phút đầu, nhiệt lượng dùng
để làm tan nước đá là:
\[
Q_{\mathrm{nc}}
=84000\cdot\frac{50}{10}
=420000\,\mathrm J.
\]

Khối lượng nước đá ban đầu:
\[
m_{\text{đ}}
=\frac{Q_{\mathrm{nc}}}{\lambda}
=\frac{420000}{3,4\cdot10^5}
\approx1,235\,\mathrm{kg}.
\]

Do đó khối lượng nước ban đầu:
\[
m_{\mathrm n}=10-1,235\approx8,765\,\mathrm{kg}.
\]

Sau $20$ phút, khối lượng nước đá đã tan:
\[
m_{\mathrm{tan}}
=1,235\cdot\frac{20}{50}
\approx0,494\,\mathrm{kg}.
\]

Khối lượng nước đá còn lại:
\[
m_{\text{còn}}
=1,235-0,494
\approx0,741\,\mathrm{kg}.
\]

Vậy:
\[
\boxed{\text{a) Đúng;\quad b) Sai;\quad c) Đúng;\quad d) Đúng.}}
\]
}
\end{ex}


% =========================================================
% BÀI 2
% =========================================================
\begin{ex}
Sự biến thiên nhiệt độ của một cục nước đá đựng trong một ca nhôm
theo nhiệt lượng cung cấp cho hệ được mô tả bằng đồ thị. Lấy nhiệt
nóng chảy riêng của nước đá là
$\lambda=3,4\cdot10^5\,\mathrm{J/kg}$, nhiệt dung riêng của nước là
$c_{\mathrm n}=4200\,\mathrm{J/(kg.K)}$, nhiệt dung riêng của nhôm là
$c_{\mathrm{Al}}=880\,\mathrm{J/(kg.K)}$. Xem hệ ca nhôm và nước luôn
có nhiệt độ giống nhau.

\begin{center}
\begin{tikzpicture}[>=stealth,scale=0.9]
    \draw[->] (0,0) -- (0,3.2) node[left] {$t(^\circ\mathrm C)$};
    \draw[->] (0,0) -- (5.5,0) node[below] {$Q(\mathrm{kJ})$};

    \draw[dashed] (0,2.5) -- (4.5,2.5);
    \draw[dashed] (3.8,0) -- (3.8,2.5);
    \draw[dashed] (4.5,0) -- (4.5,2.5);

    \draw[thick] (3.8,0) -- (4.5,2.5);

    \node[left] at (0,2.5) {$3$};
    \node[below] at (3.8,0) {$170$};
    \node[below] at (4.5,0) {$177$};

    \fill (3.8,0) circle (1.5pt);
    \fill (4.5,2.5) circle (1.5pt);

    \node[below left] at (3.8,0) {$A$};
    \node[above right] at (4.5,2.5) {$B$};
\end{tikzpicture}
\end{center}

Xét tính đúng -- sai của các phát biểu sau:

\choiceTF
{Trong giai đoạn $AB$, ca nhôm không thu nhiệt.}
{\True Cần cung cấp $170\,\mathrm{kJ}$ để làm tan chảy hoàn toàn lượng nước đá ban đầu.}
{Khối lượng nước đá ban đầu là $0,20\,\mathrm{kg}$.}
{\True Khối lượng ca nhôm xấp xỉ $0,265\,\mathrm{kg}$.}

\loigiai{
Tại điểm $A$, nước đá đã tan hoàn toàn và hệ bắt đầu tăng nhiệt độ
từ $0^\circ\mathrm C$.

Nhiệt lượng cần để làm tan hoàn toàn nước đá:
\[
Q_{\mathrm{nc}}=170\,\mathrm{kJ}=170000\,\mathrm J.
\]

Do đó khối lượng nước đá:
\[
m_{\text{đ}}
=\frac{Q_{\mathrm{nc}}}{\lambda}
=\frac{170000}{3,4\cdot10^5}
=0,50\,\mathrm{kg}.
\]

Vì vậy phát biểu c) là sai.

Từ $A$ đến $B$, nhiệt lượng cung cấp thêm:
\[
\Delta Q=177-170=7\,\mathrm{kJ}.
\]

Trong giai đoạn này, nhiệt độ tăng từ $0^\circ\mathrm C$ lên
$3^\circ\mathrm C$. Khi đó nước và ca nhôm đều thu nhiệt:
\[
\Delta Q=(m_{\mathrm n}c_{\mathrm n}+m_{\mathrm{Al}}c_{\mathrm{Al}})\Delta t.
\]

Với $m_{\mathrm n}=0,5\,\mathrm{kg}$:
\[
7000=(0,5\cdot4200+m_{\mathrm{Al}}\cdot880)\cdot3.
\]

Suy ra:
\[
m_{\mathrm{Al}}
=\frac{7000-6300}{880\cdot3}
\approx0,265\,\mathrm{kg}.
\]

Vậy:
\[
\boxed{\text{a) Sai;\quad b) Đúng;\quad c) Sai;\quad d) Đúng.}}
\]
}
\end{ex}


% =========================================================
% BÀI 3
% =========================================================
\begin{ex}
Người ta dùng một lò hồ quang điện để nấu chảy một khối kim loại
nặng $29\,\mathrm{kg}$. Biết mỗi phút, lò hồ quang cung cấp cho khối
kim loại một nhiệt lượng không đổi bằng $400\,\mathrm{kJ}$.
Sự thay đổi nhiệt độ của khối kim loại theo thời gian được biểu diễn
như hình vẽ.

\begin{center}
\begin{tikzpicture}[>=stealth,scale=0.9]
    \draw[->] (0,0) -- (0,4) node[left] {$t(^\circ\mathrm C)$};
    \draw[->] (0,0) -- (6,0) node[below] {$\tau(\text{phút})$};

    \draw[dashed] (0,3) -- (4.0,3);
    \draw[dashed] (3.0,0) -- (3.0,3);
    \draw[dashed] (4.0,0) -- (4.0,3);

    \draw[thick] (0,0.7) -- (3.0,3.0);
    \draw[thick] (3.0,3.0) -- (4.0,3.0);
    \draw[thick] (4.0,3.0) -- (5.3,3.7);

    \node[left] at (0,0.7) {$30$};
    \node[left] at (0,3) {$1530$};
    \node[below] at (3.0,0) {$50$};
    \node[below] at (4.0,0) {$70$};

    \fill (0,0.7) circle (1.5pt);
    \fill (3.0,3) circle (1.5pt);
    \fill (4.0,3) circle (1.5pt);

    \node[below right] at (0,0.7) {$A$};
    \node[above] at (3.0,3) {$B$};
    \node[above] at (4.0,3) {$C$};
    \node[above right] at (5.3,3.7) {$D$};
\end{tikzpicture}
\end{center}

Xét tính đúng -- sai của các phát biểu sau:

\choiceTF
{Giai đoạn $AB$ trên đồ thị tương ứng với quá trình nóng chảy
của kim loại.}
{Trong giai đoạn $BC$, khối kim loại không nhận thêm nhiệt lượng
từ lò nung.}
{\True Đoạn $CD$ ứng với giai đoạn khối kim loại ở thể lỏng.}
{\True Nhiệt dung riêng của khối kim loại xấp xỉ
$459,8\,\mathrm{J/(kg.K)}$.}

\loigiai{
Trong giai đoạn $AB$, nhiệt độ của kim loại tăng từ
$30^\circ\mathrm C$ lên $1530^\circ\mathrm C$, nên đây là quá trình
kim loại rắn được đun nóng đến nhiệt độ nóng chảy.

Nhiệt lượng cung cấp trong giai đoạn $AB$:
\[
Q_{AB}=400\cdot50=20000\,\mathrm{kJ}
=2,0\cdot10^7\,\mathrm J.
\]

Độ tăng nhiệt độ:
\[
\Delta t=1530-30=1500^\circ\mathrm C.
\]

Từ
\[
Q_{AB}=mc\Delta t
\]
suy ra:
\[
c=\frac{Q_{AB}}{m\Delta t}
=\frac{2,0\cdot10^7}{29\cdot1500}
\approx459,8\,\mathrm{J/(kg.K)}.
\]

Trong giai đoạn $BC$, nhiệt độ không đổi ở $1530^\circ\mathrm C$
nhưng lò vẫn cung cấp nhiệt lượng:
\[
Q_{BC}=400(70-50)=8000\,\mathrm{kJ}.
\]

Nhiệt lượng này dùng để làm nóng chảy kim loại.

Sau khi nóng chảy hoàn toàn, giai đoạn $CD$ có nhiệt độ tiếp tục
tăng nên kim loại ở thể lỏng.

Vậy:
\[
\boxed{\text{a) Sai;\quad b) Sai;\quad c) Đúng;\quad d) Đúng.}}
\]
}
\end{ex}
